\documentclass[12pt]{article}
\usepackage[utf8]{inputenc}
\usepackage[T1]{fontenc}
\usepackage[margin=1in]{geometry}
\usepackage{setspace}
\usepackage{xcolor}
\usepackage[colorlinks=true, urlcolor=blue!60!black, linkcolor=blue!60!black]{hyperref}
\setlength{\parindent}{0em}
\setlength{\parskip}{0.6em}
\renewcommand{\baselinestretch}{1.0}
\begin{document}
\begin{center}
{\LARGE \textit{Fulton v. City of Philadelphia}, 593 U.S. 522 (2021)\par}
\vspace{0.5em}
{\large Case Brief --- Two-Page Edition\par}
\end{center}

\section{Memory Jogger}
Philadelphia could not drop Catholic Social Services for refusing to certify same-sex foster couples: a discretionary-exemption clause made its policy not generally applicable under \textit{Smith} (\href{https://www.courtlistener.com/opinion/4892570/fulton-v-philadelphia/}{CourtListener}).

\section{Detailed Case Facts}
CSS, a Catholic agency that had certified foster families for the City for over 50 years, would not certify same-sex married couples; none had ever applied, and CSS would refer them elsewhere. After a news story, the City froze referrals and refused to renew CSS's contract, citing contract section 3.21, which allowed exceptions in the Commissioner's ``sole discretion,'' and its public-accommodations ordinance.

\section{Procedural History}
The district court denied a preliminary injunction under \textit{Smith}; the Third Circuit affirmed; the Supreme Court reversed.

\section{Judicial Votes}
Reversed, 9--0. Roberts wrote for six; Barrett (with Kavanaugh and, in part, Breyer) concurred; Alito and Gorsuch (each joined by Thomas and the other) concurred in the judgment.

\section{Holding}
A formal mechanism for individualized exemptions makes a policy not generally applicable even if never used, triggering strict scrutiny. Foster certification is not a public accommodation, and the City had no compelling interest in denying CSS an exception. The Court did not decide whether to overrule \textit{Smith}.

\section{Analysis of Opinions}
\textbf{Majority:} The exemption clause invites officials to judge reasons for noncompliance; strict scrutiny asks about the interest in denying \emph{this} claimant an exception, and the City's interests failed. \textbf{Barrett:} the text casts doubt on \textit{Smith}, but its replacement raises hard questions. \textbf{Alito:} overrule \textit{Smith}; the City can nullify the ruling by deleting the clause. \textbf{Gorsuch:} the majority was ``dodging'' the question.

\section{Examples: Future Applications}
\textbf{Same-Side:} a shelter code with a never-used ``good cause'' waiver triggers strict scrutiny. \textbf{Opposite-Side:} an exception-free contract or hotel-and-restaurant accommodations law stays under \textit{Smith}. \textbf{Fence-Sitter:} an exception-free rule adopted amid council members' mockery of Catholic doctrine.

\section{Critique}
Originalists say the Court ``kicked the can down the road'' by contorting the contract to avoid \textit{Smith}. Lupu, Tuttle, and Tebbe argue that a never-exercised boilerplate clause should not hollow out \textit{Smith}, and critics stress harms to LGBTQ+ foster families.

\section{Key Quotations}
\begin{quote}
``The creation of a formal mechanism for granting exceptions renders a policy not generally applicable, regardless whether any exceptions have been given, because it `invite[s]' the government to decide which reasons for not complying with the policy are worthy of solicitude, \textit{Smith}, 494 U.S. at 884---here, at the Commissioner's `sole discretion.'\,''
\end{quote}
\noindent --- Chief Justice Roberts, Opinion of the Court.

\end{document}
